\documentclass[10pt, a4paper]{article}
\usepackage[utf8]{inputenc}
\usepackage[T1]{fontenc}
\usepackage[ngerman]{babel}
\usepackage[german=swiss]{csquotes}
\usepackage{enumitem} % Für saubere Formatierung der Listen
\usepackage[left=2.5cm, right=2.5cm, top=2.5cm, bottom=2.5cm]{geometry} % Ränder etwas breiter für bessere Lesbarkeit im Einspalten-Layout

% Verhindert das Einrücken neuer Absätze für eine saubere Darstellung
\setlength{\parindent}{0pt}
\setlength{\parskip}{0.5em}

% Verhindert unschöne Seitenumbrüche (einzelne Zeilen am Anfang/Ende einer Seite)
\widowpenalty=10000
\clubpenalty=10000
\displaywidowpenalty=10000

\begin{document}

% --- TITELBLATT ---
\begin{titlepage}
    \centering
    \vspace*{4cm}
    {\Huge \textbf{Expert Briefing Document}\par}
    \vspace{1.5cm}
    {\Large Strukturvalidierung eines systemdynamischen Modells zur Landnutzungsdynamik in den EU27 (2026–2050)\par}
    
    \vspace{4cm}
    % PLATZHALTER FÜR AUTOREN
    {\large \textbf{Autoren:}\\
    Vorname Nachname\\
    Vorname Nachname\\
    Vorname Nachname\par}
    
    \vfill
    {\large \today\par}
\end{titlepage}

\newpage

% --- HAUPTDOKUMENT ---
\small

\section*{1. Zielsetzung der Anfrage}
Im Rahmen eines systemdynamischen Modellierungsprojekts untersuchen wir die strukturelle Wechselwirkung zwischen Urbanisierungsprozessen und landwirtschaftlicher Flächennutzung in den EU27 im Zeitraum 2026–2050.

Ziel dieser Anfrage ist \textbf{nicht} die Validierung von Simulationsergebnissen oder Parametrisierungen, sondern die \textbf{Strukturvalidierung der angenommenen kausalen Mechanismen} auf konzeptioneller Ebene (Causal Loop Diagram).

Wir möchten prüfen,
\begin{itemize}[itemsep=0pt, leftmargin=*]
    \item ob die modellierten Einflussbeziehungen fachlich plausibel sind,
    \item ob relevante Mechanismen fehlen,
    \item ob die angenommene Wirkungsrichtung konsistent mit dem Stand der Literatur ist,
    \item und ob zentrale Rückkopplungsstrukturen adäquat abgebildet sind.
\end{itemize}

\section*{2. Modellrahmen und Systemgrenzen}
Das Modell aggregiert die EU27 als makroskopisches Gesamtsystem. Regionale Differenzierungen werden in dieser Phase nicht modelliert.

Der Simulationszeitraum umfasst 2026–2050. Historische Daten (2000–2025) dienen der konzeptionellen Kalibrierung.

Nicht modelliert werden:
\begin{itemize}[itemsep=0pt, leftmargin=*]
    \item explizite Klimadynamiken
    \item geopolitische Schocks
    \item globale Marktgleichgewichte
    \item detaillierte sektorale Differenzierungen
\end{itemize}

Die Modellierung verfolgt einen strukturellen Ansatz im Sinne der System Dynamics (Forrester), mit Fokus auf endogene Rückkopplungen.

\section*{3. Konzeptionelle Kernmechanismen}
Auf Basis des Causal Loop Diagrams identifizieren wir vier zentrale strukturelle Dynamiken:

\subsection*{3.1 Wirtschaftlicher Verstärkungsmechanismus}
Wirtschaftliches Wachstum erhöht Einkommen, was den Flächenverbrauch pro Kopf steigert. Dies führt zu erhöhter Landumwandlung und reduziert landwirtschaftliche Fläche. Dieser Mechanismus wirkt selbstverstärkend (Reinforcing Loop), sofern keine starken Gegenmechanismen greifen.

\subsection*{3.2 Technologischer Kompensationsmechanismus}
Agrartechnologischer Fortschritt erhöht den Ertrag pro Hektar, wodurch bei gegebener Produktion weniger Fläche benötigt wird. Dieser Mechanismus wirkt potenziell stabilisierend (Balancing Loop), kann jedoch Flächenverluste nur partiell kompensieren.

\subsection*{3.3 Raumplanerischer Schutzmechanismus}
Institutionelle Eingriffe (Spatial Planning) stärken den Bodenschutz und stabilisieren landwirtschaftliche Fläche. Dieser Mechanismus ist politisch sensitiv und nicht automatisch systemisch dominant.

\subsection*{3.4 Importdruck-Mechanismus}
Sinkende Nahrungsmittelproduktion erhöht Importabhängigkeit. Dies erzeugt wirtschaftlichen Druck, der wiederum Flächennutzungsentscheidungen beeinflussen kann. Dieser Mechanismus stellt eine indirekte Rückkopplung dar.

\section*{4. Zentrale Annahmen}
\begin{itemize}[itemsep=0pt, leftmargin=*]
    \item Aggregation über EU27 ist strukturell zulässig.
    \item Land Conversion beschreibt primär die Umwandlung landwirtschaftlicher Fläche.
    \item Bevölkerung wirkt exogen auf Wohnraumnachfrage.
    \item Beziehungen werden als monoton gerichtet modelliert.
    \item Rückumwandlungsprozesse werden nicht als dominanter Trend angenommen.
\end{itemize}
Diese Annahmen sind explizit dokumentiert und offen zur Diskussion.

\section*{5. Konkrete Fragen an die Expert:innen}
\begin{enumerate}[itemsep=2pt, leftmargin=*]
    \item Sind die identifizierten Kernmechanismen relevant und korrekt?
    \item Sind die identifizierten Kernmechanismen vollständig oder fehlen strukturell relevante Einflussfaktoren?
    \item Sind die angenommenen Wirkungsrichtungen fachlich konsistent?
    \item Gibt es bekannte Nichtlinearitäten oder Schwellenwerte, die strukturell berücksichtigt werden sollten?
    \item Sind relevante zeitliche Verzögerungen konzeptionell falsch eingeschätzt?
    \item Gibt es aus Ihrer Sicht weitere übergeordnete Aspekte oder alternative strukturelle Erklärungsansätze, die wir bei der Validierung prüfen sollten?
\end{enumerate}

\section*{6. Beilage}
\begin{itemize}[itemsep=0pt, leftmargin=*]
    \item Causal Loop Diagram (CLD)
    \item Strukturierte Dokumentation der einzelnen Beziehungen
    \item Loop-Beschreibungen
\end{itemize}

\end{document}
